\scanpage{30}
\diagram{\begin{tikzcd}[ampersand replacement=\&,column sep=large,execute at end picture={
\draw ([yshift=1pt]fromgroup.south west) -- ([yshift=-2pt]fromgroup.south west) -- ([yshift=-2pt]fromgroup.south east) -- ([yshift=1pt]fromgroup.south east);
\node[below=3pt,align=center,font=\footnotesize] at (fromgroup.south) {Groupe fond.\\motivique};}]
|[alias=fromgroup]| \unsure{\Pi|\pi}(\mathrm{Spec}\,k,\unsure{\xi|\zeta})\arrow[r] \& G
\end{tikzcd}}

L'image \deleted{sur la} de la composante connexe doit être alors le sous-groupe rationnel sur $\mathbb Q$ de $G$ engendré par $i_1,i_2\colon\mathbb G_{m\,\mathbb C}\longrightarrow G_{\mathbb C}$\unsure{.} \deleted{\unsure{associé|associés}}

\uline{NB} : La seule condition 1) permet de définir un élément
\[
\psi\in Q(\mathbb C)
\]
par la condition que la bigraduation correspondante sur $T_{\mathrm{DR}}\otimes_k\mathbb C\overset{\unsure{\varphi|\psi}}{\simeq}T_B\otimes_{\mathbb Q}\mathbb C$ \unsure{satisfasse|satisfassent} à
\[
\overline{T^{pq}}=T^{q,p}.
\]

\uline{Re NB} Si \deleted{\unsure{le point}} la situation \inserted{a) à g)} provient d'une catégorie de motifs $\unsure{\mathcal M|M}$ sur $k$, alors on a dans $P\overset{G}{\times}\mathbb G_m\simeq\mathbb G_m\overset{G'}{\times}P$ un point marqué $\alpha$, qui définit donc un point $\alpha_{\mathbb C}$ de $P_{\mathbb C}\overset{G_{\mathbb C}}{\unsure{\otimes|\times}}\mathbb G_{m\,\mathbb C}$. Ceci dit, considérons $\dot{\varphi}\in(P\overset{G}{\times}\mathbb G_m)(\mathbb C)$, on aura $\dot{\varphi}=\frac{1}{2i\pi}\alpha_{\mathbb C}$ (à moins que ce soit l'inverse...), ce qui est une restriction \unsure{extrinsèque|intrinsèque} sur les données a) à g), savoir que $2i\pi\dot{\varphi}$ est $\unsure{\mathrm{rat}/k|\mathrm{rat}/k^{\smash{\prime}}}$ ....

\uline{Remarque} les données \unsure{d)|a)} à g), avec condition 1, définissent les \deleted{\unsure{$M$ etc.}} $\otimes$-catégories sur $\mathbb Q$, \inserted{de type fini, \deleted{\unsure{gradu}} \unsure{par th. $i$ et $\varepsilon$|par th. $i$ et $s$},} munies d'un $\otimes$ foncteur \unsure{comp. aux $\otimes$ et $\varepsilon$|comp. aux $\otimes$ et $\epsilon$} sur la catégorie des structures de Hodge \inserted{\unsure{enrichies|mixtes|enrichis}} $(V_{\mathbb Q},V_{\mathbb C}=\coprod V^{pq}\unsure{)|,},V_k)$, i.e. avec une \deleted{définition} \inserted{\unsure{restriction|ventilation}} \deleted{\unsure{sur}} à $k$ de \deleted{la filtration naturelle de} $V=V_{\mathbb Q}\otimes_{\mathbb Q}\mathbb C=\coprod_{pq}V^{pq}$ avec sa filtration naturelle.
% paper watermark, sideways
{\centering \textsc{Extra Stro\unsure{n}}\\
\textit{Ma\unsure{in|nn}\unsure{ol|on}\unsure{uy|ny|y} 1808}\par}
% archive sheet number: 29
